\documentclass[12pt,a4paper]{article}

\usepackage[utf8]{inputenc}
\usepackage[T2A]{fontenc}
\usepackage[russian]{babel}
\usepackage{amsmath, amssymb}
\usepackage{graphicx}
\usepackage{hyperref}
\usepackage{geometry}
\geometry{left=25mm,right=25mm,top=25mm,bottom=25mm}

\begin{document}

\title{Представление чисел в ЭВМ}
\author{Автор: Dlzxn Dev}
\date{\today}
\maketitle

\begin{abstract}
Здесь можно написать краткое резюме или аннотацию документа.
\end{abstract}

\section{Представление 0 и 1}

    0 и 1 представлены в некоторой ячейке как определенное напряжение, например,
    если напряжение >3, то это 1, если <1, то это 0. Важно заметить, что напряжение
    между этими двумя значениями не относится ни к 0, ни к 1 \
\section{Бинарные числа и логические операции}
Мы используем десятичную систему счисления, т е числа 1, 2, 3, 4, 5, 6, 7, 8, 9. Т е при записи числа мы используем степень 10.
    Например, возьмем число 898, 8 это представление \(10^0\), а вот для 9 это уже \(10^1\)
    Перейдем к двоичной системы счисления, для записи ее используется нижний индекс \(10_2\)
    Попробуем перевести \(1010_2\) в десятичную систему счисления. Разобьем все это на сумму, такую что:
    \[0*2^0 + 1*2^1 + 0*2^2+1*2^3 = 10\].
    Двоичная система счисления очень неудобна,
    т к имеет очень большую длину, поэтому для сокращения записи как правило
    используют восьмеричную и шестнадцатиричную системы счисления.


\end{document}
